%%%%%%
%
% $Author: Abishek $
% $Date: 06.05.2024 $
% $Path: BA24-02-Time-Series\report\Contents\en$
% $Version: 1.0 $
%
%
% !TeX encoding = utf8
% !TeX root = Rename
% !TeX TXS-program:bibliography = txs:///bibtex

\chapter{Documentation}

\section{Given documentation}

\subsection{Use of Best Practice document}

In our project, we diligently applied Prof. Dr. E. Wings' Scientific Project-best practices document for effective knowledge discovery and project management. We began by establishing a structured project directory and management system, following precise naming conventions for files such as "author.xlsx" and "README.md". This meticulous setup ensured clarity and organization from the outset, essential for maintaining coherence throughout our work. Additionally, we prioritized the installation and configuration of critical software tools like TeXLive, TeXstudio, Sumatra PDF, and doxygen. These tools were chosen based on their suitability for our documentation and programming needs, aligning with Prof. Dr. E. Wings' recommendations for seamless workflow integration.

Our project's foundation rested on a thorough literature review, meticulously documenting sources ranging from academic articles to technical books and online resources. This comprehensive approach not only validated our findings but also honored the contributions of previous research, ensuring our work was grounded in established knowledge. We translated this research into a detailed project plan, systematically outlining tasks, setting clear objectives, and assigning responsibilities. Regular updates to this plan enabled us to adapt to challenges and maintain progress, reflecting Prof. Dr. E. Wings' emphasis on structured project management.

In presenting our findings, we followed Prof. Dr. E. Wings' principles of effective communication. Our reports and presentations were structured to engage and inform, with clear introductions, well-supported arguments, and insightful conclusions. We adhered to rigorous documentation standards, meticulously labeling illustrations and tables, and providing comprehensive explanations for all figures used. This commitment to clarity and transparency ensured that our work was accessible and impactful, meeting the high standards expected in both academic and practical contexts.

\subsection{same notation}



		
		\begin{center}
		
			\begin{tabular}{|c|p{10cm}|} % c for centered, p{10cm} for a paragraph column with 10cm width
				\hline
				\textbf{Notation} & \textbf{Meaning} \\
				\hline
				KDD & Knowledge Discovery in Databases \\
				\hline
				CRISP-DM & Cross-Industry Standard Process for Data Mining \\
				\hline
				ML pipeline & Machine Learning Pipeline \\
				\hline
				LFS & Large File Storage \\
				\hline
			\end{tabular}
		\end{center}
		
\subsection{Supplements}
		
		Along with the Best Practices document,we also refereed other sources by Prof. Dr. E. Wings, the Blogs, Installation guides and research papers for different parts of this report and mentioned below:
		
		\begin{itemize}	
		\item Doxygen document compiled by Prof. Dr. E. Wings, provide a concise definition of what Doxygen is and its purpose in software development projects. Highlight its role as a documentation generator that extracts code structure and comments to produce readable documentation.
		\item Software installation section, we refereed several installation guides
		from official page and experts blog to avoid problems in the further
		development.
		\item In the KDD-Database section we used KDD process document compiled by Prof. Dr. E. Wings to get a complete understanding about the KDD process applied it to our examples.
			\end{itemize}
		\subsection{Same structure}
		
		 We have followed same structure from given documentation and images
		 also for following section in the report:
		\begin{itemize}	
		\item Packages
		\item Data Mining
		\item KDD-Database
		\item ML pipeline
		\item Evaluation
		\item Monitoring
	\end{itemize}
	\section{Improvements}
	We have done the following minor changes while referring the same
	structure of Scientific Project-best practices document and other documents mentinoed above in few chapters in the report and	this changes does not change the meaning of the original content the document and and all these changes will be listed in next section.
		\subsection{list of minor changes}
		
\begin{itemize}	
	\item Images captions are modified with proper spelling and citations.
	\item In GitHub-Die Basis we have translated few German words
	which are mentioned in Translations section of then report.
	\item Table information and and image labels translated in KDD section
		\end{itemize}
	
\subsection{Major improvements}
We have not modified so much in reference to old document however
we considered using books and research papers in addition to given
documents.

\subsection{Documentation of the major improvements}
NA
	\section{Tranlations}
		\subsection{1-1 translation}
		
			\begin{center}
			
			\begin{tabular}{|c|p{10cm}|} % c for centered, p{10cm} for a paragraph column with 10cm width
				\hline
				\textbf{n}German & \textbf{English} \\
				\hline
				Eckenverrundung & corner rounding \\
				\hline
				Account zulegen & Add an account \\
				\hline
			
			\end{tabular}
		\end{center}
\subsection{rewrite as an engineer}
\begin{itemize}	
	\item KDD process information are restructured in proper format and with
	relation to our project.
		\item we adhered to best practices to comprehend documentation thoroughly. This approach ensured that we could effectively grasp the nuances and requirements outlined in various project documents.
	\end{itemize}
	\subsection{Rewrite as an engineer with improvements}
	\begin{itemize}	
	\item In KDD related to our project, we write the theoretical concepts in
	relation of our project and also included the datasets and its details
	in it.
	 \item In Hardware Bill of Material section, we have updated the prices
	and working links along with data sheets for reference.
	 \item In Software Bill of Material section, we have updated the table with
	latest version of application or framework or library.
\end{itemize}